\documentclass[../main.tex]{subfiles}
\begin{document}

\section{Renormalisation Group}
\label{sec:rg}

Our main object of interest is the single-site distribution, written with the conjugate rotated, $i \hat \Phi \to \hat \Phi$, so that $\hat \Phi$ in this section denotes the (real) saddle-point value of $i \hat \Phi$ in \eqref{eq:single-site}:
\begin{align*}
  \mathcal W[\Phi, \hat \Phi]
  &= 
  \int \mathcal{D} h \exp \Bigg \{
    - \frac12 h^{\mu}_i [(\Phi + \sigma^2 I)^{-1}]^{\mu,\nu}_{i,j} h^{\nu}_j 
    + \frac12 \phi^{\mu}_i \left( \hat \Phi^{\mu,\nu}_{\alpha,\beta} K_{\alpha,i,\beta,j} - G_{i,j} \delta^{\mu,\nu}  \right) \phi^{\nu}_j
  \Bigg \}
\end{align*}
Eventually we will use the Renormalisation Group to attempt to handle this thing; but to begin with, we will develop intuition in the simplest possible case.

\subsection{Linear Activation, One Hidden Layer}
We begin by making some simplifying assumptions for the spatial kernels:
\begin{align*}
    K^{(0),\text{post}} = K^{(1),\text{pre}} \equiv K, 
    &&
    K^{(0),\text{pre}} = k_0 I,
    &&
    K^{(1),\text{post}} = k_1 I,
    &&
    G = g^2 I
\end{align*}
and studying the simplest possible case, where we have only one hidden layer. Redefining terms, we write:
\begin{align*}
    & \mathbb X \equiv \Phi^{(0)},
    \qquad
        \mathbb X^{\mu,\nu}_{i,j} =
            \frac1m
            K^{(0),\text{post}}_{i,j}
            x^{\mu}_a K^{(0),\text{pre}}_{a,b} x^{\nu}_b
        = \frac{k_0}{m} x^{\mu}_a x^{\nu}_a K_{i,j}
    \\
    & \Phi \equiv \Phi^{(1)},
    \qquad
        \Phi^{\mu,\nu}_{\alpha, \beta} =
            K^{(1),\text{post}}_{\alpha,\beta}
            \frac1N
            \left\langle
                \phi^{\mu}_i K^{(1),\text{pre}}_{i,j} \phi^{\nu}_j
            \right\rangle
        = k_1 \delta_{\alpha,\beta} \underbrace{ \frac1N \left \langle \phi^{\mu}_i K_{i,j} \phi^{\nu}_j \right \rangle }_{C^{\mu,\nu}}
\end{align*}
$C$ here \typo{measures} the aggregate correlation in the hidden layer, and is the natural analogue of the more \typo{familiar} non-spatial kernels.
For convenience we define the regularised kernel
\begin{align*}
    \tilde \Phi^{\mu,\nu}_{\alpha,\beta} 
    =
    \Phi^{\mu,\nu}_{\alpha,\beta} + \kappa^{-1} \delta^{\mu,\nu} \delta_{\alpha,\beta}
    = 
    k_1 \delta_{\alpha,\beta} \underbrace{ \left( C^{\mu,\nu} + \frac{1}{k_1 \kappa} \delta^{\mu,\nu} \right) }_{\tilde C^{\mu,\nu}}
    &&
    [\tilde \Phi^{-1}]^{\mu,\nu}_{\alpha,\beta} =  k_1^{-1} \delta_{\alpha,\beta} [\tilde C^{-1}]^{\mu,\nu}
\end{align*}
The saddle point condition is
\begin{align*}
    \hat \Phi 
    & 
    \equiv 
    \hat \Phi^{(1)} 
    =
    \tilde \Phi ^{-1} \mathbb Y \tilde \Phi^{-1}
\end{align*}
In the path integral, this $\hat \Phi^{\mu,\nu}_{\alpha,\beta}$ term is contracted with $K^{(1),\text{post}}_{\alpha,\beta} = k_1 \delta_{\alpha,\beta}$, which gives
\begin{align*}
    B^{\mu,\nu}
    \equiv
    K^{(1),\text{post}}_{\alpha,\beta} \hat \Phi^{\mu,\nu}_{\alpha,\beta}
    =
    k^{-1}_1 \delta_{\alpha,\beta} [(\tilde C)^{-1}]^{\mu,m} \delta_{\alpha,a} y^m_a y^n_b \delta_{b, \beta} [(\tilde C^{-1})]^{n,\nu}
    =
    k_1^{-1} [(\tilde C)^{-1}]^{\mu,m} y^m_{\alpha} y^n_{\alpha}  \typo{[(\tilde C)^{-1}]^{n,\nu}}
\end{align*}
or in matrix notation
\begin{align*}
    B = \frac{1}{k_1} \tilde C^{-1} Y Y^{\top} \tilde C^{-1}
\end{align*}
It is also convenient to define
\begin{align*}
    A = \frac{k_0}{m} X X^{\top}
\end{align*}

\paragraph{Diagonalising Space}

The exponent in the path integral features a kernel of the form
\begin{align*}
    ( \mathbb X^{\mu,\nu}_{i,j} + \sigma^2 \delta^{\mu,\nu} \delta_{i,j} )^{-1}
    =
    ( A^{\mu,\nu} K_{i,j} + \sigma^2 \delta^{\mu,\nu} \delta_{i,j})^{-1}
\end{align*}
We can diagonalise this in space by choosing a basis for $K$; in this case, we will use the Fourier basis (tantamount to assuming translation invariance and periodic boundaries in the $i,j$ index). We denote by $F_{i,q}$ the discrete Fourier transform matrix with plane waves of frequency $q$ in the columns, such that $K = F \Lambda F^{*}$. We can therefore write the kernel as
\begin{align*}
    A^{\mu,\nu} \underbrace{ F_{i,q} \lambda_{q} F^*_{j,q} }_{K_{i,j}} + \typo{\sigma^2} \delta^{\mu,\nu} \underbrace{ F_{i,q} F^*_{j,q} }_{\delta_{i,j}}
    =
    F_{i,q} ( \lambda_{q} A^{\mu,\nu} + \sigma^2 \delta^{\mu,\nu}) \typo{F^*_{j,q}}
\end{align*}
We therefore define the kernel that is block-diagonal in Fourier space
\begin{align*}
    \tilde A^{\mu,\nu}_{q} = \lambda_q A^{\mu,\nu} + \sigma^2 \delta^{\mu,\nu}
\end{align*}
where $\lambda_q$ are the eigenvalues of $K$. This allows us to complete the inverse
\begin{align*}
    (F_{i,q} \tilde A^{\mu,\nu}_{q} F^*_{j,q})^{-1}
    =
    F_{i,q} [(\tilde A_{q})^{-1}]^{\mu,\nu} F^*_{j,q}
\end{align*}
Similarly, we can diagonalise
\begin{align*}
    \hat \Phi^{\mu,\nu}_{\alpha,\beta} K^{(1),\text{post}}_{\alpha,\beta} K^{(1),\text{pre}}_{i,j} - g^2 \delta^{\mu,\nu} \delta_{i,j}
    =
    B^{\mu,\nu} K_{i,j} - g^2 \delta^{\mu,\nu} \delta_{i,j}
    =
    F_{i,q} ( \lambda_q B^{\mu,\nu} - g^2 \delta^{\mu,\nu} ) F^*_{j,q}
\end{align*}
by defining
\begin{align*}
    \tilde B^{\mu,\nu}_q = \lambda_q B^{\mu,\nu} - g^2 \delta^{\mu,\nu}
\end{align*}
Because the activation function is linear, $h \equiv \phi$ and we write
\begin{align*}
  \mathcal W[\Phi, \hat \Phi]
  &= 
  \int \mathcal{D} \phi \exp \Bigg \{
    - \frac12 \phi^{\mu}_i 
         F_{i,q}
         \Big(
            \underbrace{ 
                [\tilde A_q^{-1}]^{\mu,\nu}
                - \tilde B^{\mu,\nu}_q
            }_{\Sigma^{-1}_q}
        \Big)
        F^*_{j,q}
    \phi^{\nu}_j 
  \Bigg \}
\end{align*}
Without Jacobian penalty ($|F|=1$), we make the change of variables $\phi^{\mu}_i = F_{i,q} \tilde \phi^{\mu}_q \Leftrightarrow \tilde \phi^{\mu}_q = F^{*}_{i,q} \phi^{\mu}_i$
\begin{align*}
  \mathcal W[\Phi, \hat \Phi]
  &= 
  \int \mathcal{D} \tilde \phi \exp \Bigg \{
    - \frac12 \tilde \phi^{\mu}_q 
         [\Sigma^{-1}_q]^{\mu,\nu}
    \tilde \phi^{\nu *}_q 
  \Bigg \}
\end{align*}
We have here obtained the frequency-dependent data-interaction matrix $\Sigma_q$. For each $q$, this matrix can in principle have a different eigenbasis--these are the directions in data-space that define our functional cell types: the eigendirections of a low-frequency data interaction matrix describe response profiles common to all cells; the directions associated with an intermediate frequency describe functional variations between, say, areas, or equivalently describe \typo{coarse} functional variation; and the directions associated with high frequency modes describe functional cells. So the task then is to \typo{diagonalise} these $\Sigma_q$ and ask how these directions vary with $q$.

\paragraph{Diagonalising Function}

The goal now is to diagonalise
\begin{align*}
    \Sigma_q = \left( \tilde A_q^{-1} - \tilde B_q \right)^{-1} 
    &&
    \text{with}
    &&
    \tilde A_q = \lambda_q \frac{k_0}{m} X X^{\top} + \sigma^2 I_P
    &&
    \text{and}
    &&
    \tilde B_q = \frac{\lambda_q}{k_1} \tilde C^{-1} Y Y^{\top} \tilde C^{-1} - g^2 I_P
\end{align*}
Consider firstly the $C$ kernel lurking inside the $B$ term:
\begin{align*}
    C^{\mu,\nu}
    =
    \frac1N \sum_{i,j} K_{i,j} \left \langle \phi^{\mu}_i \phi^{\nu}_j \right \rangle
    =
    \frac1N \sum_{i,j,q,q'} K_{i,j} F_{i,q} \left \langle \tilde \phi^{\mu}_q \tilde \phi^{\nu}_{q'} \right \rangle F_{j, q'}
    =
    \frac1N \sum_{i,j,q,q'} K_{i,j} F_{i,q} \delta_{q,-q'} \Sigma^{\mu,\nu}_q F_{j, q'}
    =
    \frac1N \sum_q \lambda_q \Sigma^{\mu,\nu}_q
\end{align*}
where the $\sum_q \lambda_q$ appeared from
\begin{align*}
    \sum_{i,j,q'} K_{i,j} F_{i,q} \delta_{q,-q'} F_{j,q'} 
    =
    \sum_{i,j} K_{i,j} F_{i,q} F^*_{j,q}
    =
    \sum_{i,j,q'} F_{j,q'} \lambda_{q'} F^*_{i,q'} F_{i,q} F^*_{j,q}
    =
    \sum_{q'} \delta_{q,q'} \lambda_{q'}
    =
    \lambda_q
\end{align*}
That $C$, and therefore $B$, is defined in terms of $\Sigma$, exposes the fundamental difficulty that if we are to diagonalise any one $\Sigma_q$, we have to make it consistent with those of all the $q$s.
We therefore make the ansatz
\begin{align*}
    \Sigma_q = U D_q U^{\top} \quad \forall q
\end{align*}
which is to say that all spatial frequencies express function using the same basis, $U$; the extent to which functional variation manifests on different spatial scales is captured by $D_q$.
To make progress, we will also assume that this basis $U$ diagonalises the data kernels $XX^{\top}$ and $Y Y^{\top}$, which is possible only if $[XX^{\top}, YY^{\top}] = 0$:
\begin{align*}
    X = U S_X V_X^{\top}
    \implies
    XX^{\top} = U S_X^2 U^{\top}
    &&
    Y = U S_Y V_Y^{\top}
    \implies
    YY^{\top} = U S_Y^2 U^{\top}
\end{align*}
This ansatz immediately throws out the possibility of answering one of our motivating questions, which was how data-directions vary with spatial scale (now all spatial scales express function in the same vocabulary, the columns of $U$). It is, however, manifestly self-consistent: if we start from a $\Sigma_q$ so diagonalised, and iterate to a fixed point, it will stay so diagonlised. For generic data this assumption of mutual diagonalisability is an approximation that we should like to examine more deeply.
Now we can diagonalise $\tilde A$:
\begin{align*}
    \tilde A_q 
    = 
    U \Big( \lambda_q \frac{k_0}{m} S_X^2 + \sigma^2 I_P  \Big) U^{\top}
    \equiv
    U D_{A,q} U^{\top}
\end{align*}
Under the ansatz, we can also diagonalise $\tilde C$ and $\tilde B$:
\begin{align*}
    \tilde C
    =
    \frac1N \sum_q \lambda_q \Sigma_q + \frac{1}{k_1 \kappa} I_P
    =
    U \Big( \frac1N \sum_q \lambda_q D_q + \frac{1}{k_1 \kappa} I_P \Big) U^{\top}
    \equiv
    U D_{C} U^{\top}
\end{align*}
and consequently
\begin{align*}
    \tilde B_q 
    = 
    U \Big( \frac{\lambda_q}{k_1} D_{C}^{-2} S_Y^2 - g^2 I_P \Big) U^{\top}
    \equiv 
    U D_{B,q} U^{\top}
\end{align*}
Thus we have diagonalised the problem
\begin{align*}
    \Sigma_q
    =
    U D_q U^{\top} 
    = 
    U \Big(D_{A,q}^{-1} - D_{B,q} \Big)^{-1} U^{\top} 
\end{align*}
\begin{align*}
    D_{A,q} = \lambda_q \frac{k_0}{m} S_X^2 + \sigma^2 I_P,
    &&
    D_{B,q} = \frac{\lambda_q}{k_1} D_C^{-2} S_Y^2 - g^2 I_P,
    &&
    D_C = \frac1N \sum_q \lambda_q D_q + \frac{1}{k_1 \kappa} I_P
\end{align*}
Let us denote by $d_q$ an element on the diagonal of $D_q$; by $a_q$ and $b_q$ elements on the diagonals of $D_{A,q}$ and $D_{B,q}$, respectively; and by $c$ an element on the diagonal of $D_C$. We define
\begin{align*}
    \hat s_X^2 = \frac{k_0 s_X^2}{m}
    &&
    \hat s_Y^2 = \frac{s_Y^2}{k_1}
    &&
    \hat \sigma^2_q = \frac{\typo{\sigma^2}}{\lambda_q}
    &&
    \delta = \frac{1}{k_1 \kappa}
\end{align*}
to solve:
\begin{align*}
    d_q^{-1}
    =
    a_q^{-1} - b_q
    =
    \frac{1}{\lambda_q \hat s_X^2 + \sigma^2} - \frac{\lambda_q \hat s_Y^2}{c^2} + g^2
    =
    \frac{\lambda_q}{h_q} - \frac{\lambda_q \typo{\hat s_Y^2}}{c^2}
\end{align*}
where we have introduced
\begin{align*}
    h_q 
    =
    \frac{\lambda_q \hat s_X^2 + \sigma^2 }{\lambda_q^{-1} + g^2 (\hat s_X^2 + \hat \sigma^2_q)}
\end{align*}
this gives
\begin{align*}
    d_q = \left(\frac{\lambda_q}{h_q} - \frac{\lambda_q \hat s_Y^2}{c^2} \right)^{-1}
    =
    \frac{c^2 h_q / \lambda_q}{c^2 - \hat s_Y^2 h_q}
\end{align*}
and therefore
\begin{align*}
    c 
    = 
    \frac1N \sum_q \lambda_q d_q  + \delta
    =
    \frac{c^2}{N} \sum_q \frac{h_q}{c^2 - \hat s_Y^2 h_q} + \delta 
\end{align*}
which leaves the final implicit solution
\begin{align}
    1 - \frac{\delta}{c} = \frac{c}{N} \sum_q \frac{h_q}{c^2 - \hat s_Y^2 h_q}
    \label{eq:linear-soln}
\end{align}

\paragraph{Existence and Uniqueness}

The physical solution to \eqref{eq:linear-soln} is the one that must satisfy the positive definiteness of $\Sigma_q$, or equivalently the positivity of $d_q$, which therefore requires $c > \hat s_Y \max_q \lambda_q$.

If we define the residual of \eqref{eq:linear-soln} from its fixed point, $f(c) = 1 - \frac{\delta}{c} - \frac{c}{N} \sum_q h_q/(c^2 - \hat s_Y^2 h_q)$, then we can track how the residual changes on the physical branch:
\begin{align*}
    f'(c) 
    = 
    \frac{\delta}{c^2} 
    + 
    \frac{1}{N} \sum_q \frac{h_q (c^2 + \hat s_Y^2 h_q)}{(c^2 - \hat s_Y^2 h_q)^2}
\end{align*}
which is always positive. At the pole $c = \hat s_Y^2 \max_q h_q$, $f \rightarrow - \infty$ and increases monotonically to $1$ as $c \rightarrow \infty$, and therefore the root $f(c^{\star})=0$ is the unique root on this physical branch.

\end{document}
